\chapter{Projekt eksperymentu referencyjnego i metryki ewaluacji autoremediacji}

W celu rzetelnej oceny wpływu potoku ACE na skuteczność systemu, proces badawczy został podzielony na dwie główne fazy: fazę referencyjną (Baseline) oraz fazę ewolucyjną. Kluczowym założeniem było zachowanie determinizmu warunków początkowych poprzez wykorzystanie tej samej, losowo wygenerowanej sekwencji 100 incydentów (scenariuszy usterki) dla obu wariantów.

\section{Konfiguracja eksperymentu referencyjnego (Baseline)}

Pierwsza faza eksperymentu polegała na uruchomieniu agenta SRE w trybie „bez uczenia” (\textit{Agent No Learning}). W tym wariancie potok ACE był całkowicie wyłączony, a agenci korzystali wyłącznie z bazowych, statycznych promptów systemowych. 

Proces ten przebiegał zgodnie z logiką zawartą w module \texttt{episodes\_runner/runner.py}:
\begin{enumerate}
    \item \textbf{Przygotowanie środowiska}: Dla każdego z 100 incydentów tworzona była czysta kopia robocza aplikacji (\textit{workspace}).
    \item \textbf{Wstrzyknięcie usterki}: Aplikacja była sabotowana za pomocą wybranego patcha lub manifestu Chaos Mesh z repozytorium \texttt{fault-vault}.
    \item \textbf{Próba naprawy}: Agent SRE (\textit{Incident Commander} wraz z deputatami) podejmował próbę remediacji bez możliwości aktualizacji swojej bazy wiedzy między epizodami.
    \item \textbf{Persystencja}: Wyniki działań, ślady rozumowania i wywołania narzędzi były zapisywane w bazie \texttt{agent\_no\_learning.db}.
\end{enumerate}

\section{Eksperyment ewolucyjny z potokiem ACE}

W drugiej fazie wykorzystano ten sam zestaw 100 incydentów, jednak z włączonym mechanizmem samouczenia. Zgodnie z implementacją w \texttt{episodes\_runner/experiment\_runner.py}, potok ACE nie był uruchamiany po każdym epizodzie, lecz w interwałach co 5 pomyślnie zakończonych lub przerwanych zadań (warunek \texttt{episode\_count \% 5 == 0}).

Taka strategia (\textit{batch learning}) ma na celu stabilizację procesu uczenia i uniknięcie zbyt gwałtownych zmian w Playbookach na podstawie pojedynczych, specyficznych przypadków. Pozwala to na akumulację doświadczeń z kilku różnorodnych incydentów przed dokonaniem syntezy i konsolidacji wiedzy przez agenta Kuratora.

\section{Metryka ewaluacji: Agent Sędzia i Dowody Wykonania}

Ewaluacja każdego epizodu odbywa się w sposób zautomatyzowany przy pomocy agenta sędziego (\textit{Judge Agent}), co opisano w \texttt{experiment\_runner.py}. Sędzia dokonuje binarnej oceny (0 lub 1) w trzech kluczowych kategoriach:
\begin{itemize}
    \item \textbf{root\_cause\_analysis}: Czy agent poprawnie zidentyfikował faktyczną przyczynę awarii (zgodnie z usterką w \texttt{FAULT.md})?
    \item \textbf{successful\_fix}: Czy zastosowano poprawną poprawkę w kodzie/konfiguracji ORAZ czy wywołano narzędzie \texttt{deploy\_app}? Sędzia weryfikuje to na podstawie raportu \textit{System Execution Evidence}, ignorując deklaracje agenta, które nie znalazły odzwierciedlenia w rzeczywistych wywołaniach narzędzi.
    \item \textbf{system\_recovery\_visible}: Czy metryki z Prometheusa i Loki potwierdzają powrót systemu do stanu zdrowego (brak błędów 5XX, stabilizacja zasobów)?
\end{itemize}

\section{Cele analizy i zjawiska emergentne}

Głównym celem porównania obu faz eksperymentu, realizowanym przez skrypt \texttt{analyze\_tasks.py}, jest weryfikacja następujących hipotez:

\begin{enumerate}
    \item \textbf{Poprawa efektywności operacyjnej}: Czy wraz z ewolucją Playbooka spadnie średni czas naprawy (TTR) oraz liczba iteracji potrzebnych do rozwiązania problemu?
    \item \textbf{Redukcja błędów narzędziowych}: Czy precyzyjne instrukcje IF/THEN w Playbooku ograniczą liczbę błędnych wywołań narzędzi (\textit{tool error rate}) i halucynacji parametrów?
    \item \textbf{Emergentna organizacja pracy}: Czy system wypracuje zaawansowane wzorce delegacji? Przykładowo, czy \textit{Incident Commander} nauczy się zlecać analizę logów agentowi metryk (\textit{Monitoring Agent}) przed podjęciem decyzji o zmianie kodu przez \textit{DevOps Agenta}, co jest widoczne w strukturze drzewa zadań (\textit{tasks tree}).
\end{enumerate}

Wyniki eksperymentu są prezentowane na chronologicznych wykresach porównawczych, które umożliwiają ocenę zarówno poprawy, jak i ewentualnej regresji skuteczności agentów w poszczególnych epizodach. Na wykresach tych można obserwować momenty, w których potok ACE aktualizuje Playbook (oznaczane jako pionowe linie przerywane), co często koreluje z pojawieniem się nowych, skutecznych heurystyk działania — zwłaszcza w trudniejszych klasach incydentów. Wykresy pozwalają również wychwycić epizody, w których dochodzi do chwilowego spadku skuteczności względem fazy Baseline.

Oprócz kategorii sukcesu ocenianych przez agenta sędziego, do analizy wykorzystano także szereg metryk opisowych takich jak: czas do przywrócenia sprawności (TTR), liczba iteracji (\textit{iterations count}), częstość i odsetek błędów wywołań narzędzi (\textit{tool error count}, \textit{tool error rate}), długość trajektorii rozwiązywania zadania, liczba i głębokość potomnych zadań oraz inne statystyki syntetyczne. Statystyki te umożliwiają szczegółową ocenę wpływu ewolucji Playbooka na efektywność operacyjną systemu.
